\chapter{Stockage, transport et optionnalité physique}\label{ch:s2:storage-transport-and-physical-optionality}

Un trader qui loue une caverne de stockage de gaz détient une série d'options sur calendar spread~; celui qui détient une capacité de pipeline détient des options sur \omterm{def:s2:commodity-spreads:location}{spread géographique}. Tout
actif physique se valorise comme la somme des options qu'il contient, et le travail du trader consiste à les transformer en argent au moyen de couvertures. Secomandi a établi que l'essentiel de la
valeur de l'adaptation à l'incertitude sur les prix peut être obtenue simplement en réoptimisant, encore et encore, un plan qui ignore cette incertitude. Sur le
bail synthétique de ce chapitre, verrouiller dès aujourd'hui le meilleur plan vaut 716\,000~\$. Réoptimiser ce plan chaque mois sur la nouvelle courbe à terme et re-couvrir ajoute
56\,800~\$ en moyenne, et sur aucun chemin il ne rapporte moins. À deux centimes par unité de forward négociée, il ajoute 26\,600~\$ et perd face au plan verrouillé sur un tiers
des chemins. La construction est \texttt{firm.physopt}.

\section{Le stockage comme option}

\begin{definition}[Contrat de stockage]\label{def:s2:storage-transport-and-physical-optionality:storage}
Un \emph{contrat de stockage}\index{contrat de stockage} est un bail de capacité de stockage (volume utile et débits d'injection et de soutirage, pour une période) que le
trader monétise en achetant la matière première lorsqu'elle est bon marché, en la conservant puis en la vendant lorsqu'elle est chère, et couvre chaque injection et chaque soutirage prévus sur le
marché à terme.
\end{definition}

Le bail de \texttt{firm.physopt} porte sur 10 unités de 100\,000 MMBtu, permet d'injecter jusqu'à 2 unités et de soutirer jusqu'à 4 unités par mois, et coûte 2 centimes par MMBtu
dans chaque sens. La courbe à terme et le modèle de prix sont ceux du livre~6 (chapitre~16)~: une courbe illustrative de type Henry Hub, allant d'avril (2.60~\$) à janvier
(3.50~\$) et un modèle à deux facteurs avec une volatilité de 60~\% en bout court et de 20~\% en bout long. Le plan intrinsèque achète 2 unités chaque mois, d'avril à août
et en vend 4 en décembre, 4 en janvier et 2 en février. Verrouillé par des forwards, il vaut 716\,000~\$.

Quatre méthodes d'évaluation donnent quatre nombres~:

\begin{center}
\small
\begin{tabular}{@{}p{11.4cm}r@{}}
\toprule
méthode & valeur (milliers de dollars)\\
\midrule
intrinsèque~: meilleur plan du jour, verrouillé & 716\\
options sur spread~: les dix paires injection-soutirage du plan, chacune étant une option & 722\\
Monte Carlo par moindres carrés sur le spot (livre~6) & 760\\
intrinsèque glissant~: réoptimisé chaque mois, couvertures négociées & 773\\
\bottomrule
\end{tabular}
\end{center}

La méthode par options sur spread associe chaque unité injectée à une unité soutirée, premier entré premier sorti, et valorise chaque paire comme une option sur calendar spread portant sur les
deux forwards, de prix d'exercice égal aux coûts d'exploitation (le \cref{lst:s2:storage-transport-and-physical-optionality:options}). Ici les paires sont profondément dans la monnaie
et n'ajoutent guère. La méthode Monte Carlo par moindres carrés ne s'appuie que sur les prix spot et estime les valeurs de continuation par régression, ce qui fait perdre un peu.
L'intrinsèque glissant utilise chaque mois toute la courbe à terme. Lai, Margot et Secomandi ont montré que les heuristiques des praticiens sont rapides mais nettement
sous-optimales, et qu'avec une réoptimisation périodique au sein d'une simulation Monte Carlo elles deviennent presque optimales, à une exception près.

\section{Monétiser les options}

La valeur se réalise grâce à un programme de couverture (le \cref{lst:s2:storage-transport-and-physical-optionality:programme}). Le programme statique achète
des forwards pour chaque injection prévue et en vend pour chaque soutirage prévu, soit 20 unités au total, et livre le plan au spot. Quoi que
fassent les prix, il rapporte 716\,000~\$~: les gains et pertes des couvertures compensent exactement les prix spot payés et reçus. Le programme glissant réoptimise le
plan chaque mois sur la courbe à terme simulée, compte tenu du gaz déjà en stock, et négocie la variation des couvertures. Un nouveau plan n'est retenu que si
il vaut davantage sur la courbe courante que l'ancien, et les couvertures verrouillent ce gain~; ainsi, sans coûts de transaction, le programme glissant ne peut jamais rapporter moins que
le programme statique sur tout chemin.

\begin{center}
\small
\begin{tabular}{@{}p{5.6cm}cccc@{}}
\toprule
coût par unité de forward négociée & 0 & 0.5 centime & 2 centimes & 2 centimes, filtré\\
\midrule
surplus par rapport au programme statique (milliers de dollars) & 56.8 & 49.2 & 26.6 & 37.4\\
chemins sous le programme statique & 0~\% & 3.1~\% & 32.6~\% & 0~\%\\
unités supplémentaires de forward négociées & 15.1 & 15.1 & 15.1 & 9.2\\
\bottomrule
\end{tabular}
\end{center}

Le programme statique paie aussi ses propres 20 unités de couverture, si bien que la comparaison se fait à coûts égaux. À deux centimes par unité, le programme glissant
rapporte encore davantage en moyenne, mais il perd face au statique sur un tiers des chemins, car il re-couvre dès que le nouveau plan est meilleur, quel que soit le montant
(la \cref{fig:s2:storage-transport-and-physical-optionality:rolling}). La version filtrée ne re-couvre que lorsque le gain sur la courbe dépasse le coût des
transactions. Elle négocie 9.2 unités supplémentaires au lieu de 15.1, rapporte en moyenne 37\,400~\$ de plus que le programme statique, et jamais moins.

\begin{omfigure}
\begin{tikzpicture}
\begin{axis}[width=11cm,height=5.4cm,xlabel={\small programme glissant moins programme statique (\$ par unité de MMBtu~; 100\,000~\$ par unité)},
  ylabel={\small part des chemins (\%)},tick label style={font=\footnotesize},xmin=-1,xmax=3,ymin=0,ymax=26,grid=major,grid style={gray!25},
  table/col sep=comma,legend style={font=\scriptsize,at={(0.97,0.97)},anchor=north east}]
\addplot[const plot,omThm,thick] table[x=lo,y=free]{figdata/strategies-2/22-storage-transport-and-physical-optionality/rolling.csv};
\addplot[const plot,omDef,thick,dashed] table[x=lo,y=cost2]{figdata/strategies-2/22-storage-transport-and-physical-optionality/rolling.csv};
\legend{sans coût de transaction,2 centimes par unité}
\end{axis}
\end{tikzpicture}

\omcaption{Distribution sur 2\,000 chemins simulés du P\&L du programme en intrinsèque glissant, moins celui du programme statique, sans coût de transaction et
à 2 centimes par unité de forward négociée~; les classes extrêmes regroupent les queues. Données~: \texttt{s2\_physopt.programmes}.}
\label{fig:s2:storage-transport-and-physical-optionality:rolling}
\end{omfigure}

\section{Le transport comme option}

\begin{definition}[Option de transport]\label{def:s2:storage-transport-and-physical-optionality:transport}
Une \emph{option de transport}\index{option de transport} est le droit, conféré par une capacité de pipeline ou de transport maritime, de déplacer une matière première d'un lieu à
un autre sur une période, à un coût variable~; elle vaut une option sur spread portant sur le prix de destination moins le prix d'origine, dont le prix d'exercice est ce coût.
\end{definition}

Une capacité de 1 unité par mois depuis le hub vers un hub de zone de marché, dont les forwards dépassent ceux du hub de 5 centimes au printemps à 60 centimes en janvier, avec un
coût variable de 10 centimes, vaut 139\,000~\$ sur les forwards du jour~: seuls les mois d'hiver rapportent. En tant que strip d'options sur spread, avec des volatilités de 45~\%
et 55~\%, une corrélation de 0.85, il vaut 369\,000~\$ (la \cref{fig:s2:storage-transport-and-physical-optionality:transport}). Sa valeur dépend de la
manière dont les deux hubs évoluent indépendamment~: 580\,000~\$ pour une corrélation de 0.5 et 273\,000~\$ pour 0.95. Le détenteur le monétise comme un stockage, en vendant le
spread à terme sur les mois où il est dans la monnaie, et en réoptimisant à mesure que la base évolue.

\begin{omfigure}
\begin{tikzpicture}
\begin{axis}[width=11cm,height=5.2cm,ybar,bar width=6pt,xlabel={\small mois de livraison},ylabel={\small \$ par MMBtu},
  tick label style={font=\footnotesize},ymin=0,ymax=0.8,xtick=data,xticklabels={avr.,mai,juin,juil.,août,sept.,oct.,nov.,déc.,janv.,févr.,mars},x tick label style={rotate=40,anchor=north east},
  grid=major,grid style={gray!25},table/col sep=comma,legend style={font=\scriptsize,at={(0.03,0.97)},anchor=north west}]
\addplot[fill=gray!50,draw=gray] table[x=k,y=intrinsic]{figdata/strategies-2/22-storage-transport-and-physical-optionality/transport.csv};
\addplot[fill=omThm!60,draw=omThm] table[x=k,y=option]{figdata/strategies-2/22-storage-transport-and-physical-optionality/transport.csv};
\legend{intrinsèque,option sur spread}
\end{axis}
\end{tikzpicture}

\omcaption{Une année de capacité de transport mensuelle depuis le hub vers un hub de zone de marché~: la valeur intrinsèque sur les prix à terme d'aujourd'hui et la valeur de l'option sur spread
de chaque mois, par MMBtu. Données~: \texttt{s2\_physopt.transport\_strip}.}\label{fig:s2:storage-transport-and-physical-optionality:transport}
\end{omfigure}

Une cargaison de gaz naturel liquéfié (GNL) destinée à l'Europe à 10.00~\$ par MMBtu, qui peut à la place partir vers l'Asie à 10.50~\$, pour 80~cents de fret supplémentaire,
n'a aucune valeur intrinsèque dans le changement de destination. Ce changement reste une option sur le prix asiatique moins le prix européen, au prix d'exercice de 80~cents~: avec trois mois
pour décider, des volatilités de 60~\% et une corrélation de 0.8, elle vaut 64~cents par MMBtu, et 1.08~\$ pour une corrélation de 0.5.

\section{Ce que possède le trader adossé à des actifs physiques}

\begin{definition}[Trading adossé à des actifs physiques]\label{def:s2:storage-transport-and-physical-optionality:assetbacked}
Le \emph{trading adossé à des actifs physiques}\index{trading adossé à des actifs physiques} (asset-backed trading) consiste à trader autour d'actifs physiques que la firme possède ou loue (stockage, capacité de transport,
production, cargaisons), à les valoriser comme des options et à monétiser leur flexibilité au moyen de couvertures sur les marchés à terme et d'options.
\end{definition}

Le trader adossé à des actifs physiques possède trois éléments qu'un trader financier ne possède pas. Le premier est constitué des options elles-mêmes, qui ne peuvent pas être achetées en bourse.
Le deuxième est l'obligation de les exploiter, avec les contraintes physiques qui rendent les options précieuses et les plans compliqués. Le troisième est
l'information~: les flux qui traversent l'actif révèlent l'offre et la demande avant que les prix ne le fassent. Ce que le trader ne possède pas, c'est la certitude que les spreads seront au rendez-vous.
De 2007 à 2025, le prix spot de Henry Hub en hiver (de décembre à février) a été supérieur à celui de l'été précédent (d'avril à juin) dans 12 années sur 19. Il
était en moyenne inférieur de 7~cents, avec un écart-type de 2.04~\$. En 2008, l'hiver était inférieur de 6.12~\$ à l'été, et en 2022 de 3.71~\$. Un opérateur qui injectait
et attendait sans couvrir gagnait ces montants. Le programme de couverture existe pour les remplacer par le spread à terme au moment de la décision.

\section{Fiches de stratégie}

\begin{strategyfile}{Stockage à intrinsèque glissant}\label{strat:s2:storage-transport-and-physical-optionality:rolling}
\sfield{Qui paie, et pourquoi} Des consommateurs qui paient une prime hivernale~; le spread saisonnier de la courbe à terme rémunère le fait de détenir du gaz pendant l'été.
\sfield{Instruments et plateformes} Un \omterm{def:s2:storage-transport-and-physical-optionality:storage}{contrat de stockage}~; des contrats à terme mensuels sur gaz.
\sfield{Signal} Le meilleur plan sur la courbe à terme actuelle.
\sfield{Dimensionnement et exécution} Couvrir l'ensemble du plan~; réoptimiser chaque mois ou lors de mouvements importants~; ne recouvrir que lorsque le gain dépasse le coût.
\sfield{Coûts} Le \omterm{def:s2:storage-transport-and-physical-optionality:storage}{contrat de stockage}~; les frais d'injection et de soutirage~; la fourchette à terme.
\sfield{Comment elle échoue} Des coûts de contrat supérieurs à la valeur~; des courbes plates~; des interruptions d'exploitation.
\sfield{Horizon, capacité, infrastructure} Une année de stockage~; la programmation et les nominations.
\sfield{Backtest honnête} Courbes à terme telles qu'elles étaient~; contraintes physiques respectées~; coûts sur chaque recouvrement.
\sfield{Sources} Secomandi (2010)~; Lai, Margot et Secomandi (2010)~; ce chapitre~: 56\,800~\$ au-dessus de l'intrinsèque sans coûts.
\end{strategyfile}

\begin{strategyfile}{Couverture de stockage par option sur spread}\label{strat:s2:storage-transport-and-physical-optionality:options}
\sfield{Qui paie, et pourquoi} Des acheteurs d'options sur calendar spread qui veulent la flexibilité dont dispose le propriétaire du stockage.
\sfield{Instruments et plateformes} Des options sur calendar spread sur contrats à terme gaz~; des options sur spread de gré à gré.
\sfield{Signal} La valeur d'option du stockage comparée au prix de marché des options.
\sfield{Dimensionnement et exécution} Vendre les options que l'installation peut honorer~: des paires de mois d'injection et de soutirage.
\sfield{Coûts} La fourchette sur les options.
\sfield{Comment elle échoue} Vendre davantage d'optionalité que ne le permettent les débits de l'installation.
\sfield{Horizon, capacité, infrastructure} Une saison.
\sfield{Backtest honnête} Exercice contraint par les débits de l'installation.
\sfield{Sources} Ce chapitre~: les paires valent 722\,000~\$ contre 716\,000~\$ d'intrinsèque.
\end{strategyfile}

\begin{strategyfile}{Optimisation de la capacité de transport}\label{strat:s2:storage-transport-and-physical-optionality:transport}
\sfield{Qui paie, et pourquoi} Des acheteurs dans des zones de marché contraintes~; la base rémunère le gazoduc.
\sfield{Instruments et plateformes} Capacité de gazoduc~; swaps ou contrats à terme sur base aux deux hubs.
\sfield{Signal} La base à terme comparée au coût variable.
\sfield{Dimensionnement et exécution} Vendre la base à terme sur les mois dans la monnaie~; réoptimiser à mesure qu'elle bouge.
\sfield{Coûts} Redevances de capacité~; carburant~; fourchette sur la base.
\sfield{Comment elle échoue} De nouveaux gazoducs qui résorbent la base~; une corrélation qui augmente.
\sfield{Horizon, capacité, infrastructure} Des saisons à des années~; nominations.
\sfield{Backtest honnête} Base aux deux hubs telle qu'elle était~; interruptions de capacité.
\sfield{Sources} Ce chapitre~: 139\,000~\$ d'intrinsèque, 369\,000~\$ en options.
\end{strategyfile}

\begin{strategyfile}{Détournement de cargaison}\label{strat:s2:storage-transport-and-physical-optionality:cargo}
\sfield{Qui paie, et pourquoi} Des marchés qui enchérissent sur une offre flexible lorsqu'elle se raréfie.
\sfield{Instruments et plateformes} Cargaisons de GNL à flexibilité de destination~; contrats à terme régionaux sur gaz~; fret.
\sfield{Signal} Les netbacks par destination~: le prix moins le transport maritime et la regazéification.
\sfield{Dimensionnement et exécution} Couvrir la destination engagée~; conserver le changement de destination comme une option.
\sfield{Coûts} Transport maritime, évaporation (boil-off), redevances de canal et de port.
\sfield{Comment elle échoue} Restrictions contractuelles~; envolées du fret~; convergence des prix régionaux.
\sfield{Horizon, capacité, infrastructure} De quelques semaines à quelques mois~; transport maritime et programmation.
\sfield{Backtest honnête} Fret et netbacks à la date de décision.
\sfield{Sources} Ce chapitre~: 64~cents par MMBtu pour une corrélation de 0.8.
\end{strategyfile}

\section{Tutoriel~: l'actif comme options}\label{tut:s2:storage-transport-and-physical-optionality:tut}

\begin{tutorial}
\textbf{Objectif.} Valoriser un \omterm{def:s2:storage-transport-and-physical-optionality:storage}{contrat de stockage} de quatre façons, faire tourner ses programmes de couverture statique et glissante chemin par chemin avec les coûts de transaction, et valoriser la
capacité de transport et un changement de destination de cargaison comme des options sur spread. \textbf{État final~:} les deux tableaux et les deux figures.
\begin{tutsteps}
\item \textbf{Le programme de couverture}.
\omcode{code/firm/physopt/firm_physopt.py}{60}{91}{Verrouillage du plan, réoptimisation et recouvrement quand cela rapporte.}\label{lst:s2:storage-transport-and-physical-optionality:programme}
\item \textbf{Les options}.
\omcode{code/firm/physopt/firm_physopt.py}{94}{124}{Paires d'options sur spread, transport et détournement.}\label{lst:s2:storage-transport-and-physical-optionality:options}
\item \textbf{Exécuter} \texttt{storage()}, \texttt{hedging()}, \texttt{gated()}, \texttt{transport\_strip()}, \texttt{cargo()},
\texttt{fig\_physopt.py} et, une fois, \texttt{s2\_fetch\_hh.py}.
\end{tutsteps}
\textbf{Pour aller plus loin.} Réoptimiser uniquement sur les mouvements importants de la courbe~; ajouter le prix du \omterm{def:s2:storage-transport-and-physical-optionality:storage}{contrat de stockage} et trouver le seuil de rentabilité~; couvrir la série de transport et
mesurer son profil de P\&L chemin par chemin.
\end{tutorial}

\section{Construction~: optionalité physique}\label{bld:s2:storage-transport-and-physical-optionality:physopt}

\begin{build}
\bfield{Rôle} Programmes de couverture de stockage chemin par chemin, décomposition en options sur spread, séries de transport et changements de destination de cargaison, sur le modèle énergétique du livre~6.
\bfield{Interface} \texttt{plans(fac, prices, inv)}, \texttt{hedge\_programme(fac, model, F0, T, paths, seed, cost, roll, gate)},
\texttt{spread\_option\_pairs(fac, F0, T, model)}, \texttt{transport(FA, FB, T, cost, sA, sB, rho)},
\texttt{diversion(F1, F2, K, T, s1, s2, rho)}.
\bfield{Règles} Couvertures longues pour les injections et courtes pour les soutirages~; valorisation à la courbe de chaque mois~; physique au spot~; coûts par unité de contrat à terme échangé.
\bfield{Tests d'acceptation} \texttt{code/firm/physopt/tests/}~: les plans égalent l'intrinsèque du livre~6~; le programme statique rapporte la valeur intrinsèque sur chaque
chemin, et le programme glissant jamais moins sans coûts~; paires, transport et détournement vérifiés à la main.
\bfield{Pour aller plus loin} Rachets (débits dépendant du niveau de stock)~; valorisation de \omterm{def:s2:storage-transport-and-physical-optionality:storage}{contrats de stockage}~; contrats swing.
\end{build}

\begin{omsources}
\item N.~Secomandi, ``Optimal commodity trading with a capacitated storage asset'', \emph{Management Science} 56(3), 2010.
\item G.~Lai, F.~Margot et N.~Secomandi, ``An approximate dynamic programming approach to benchmark practice-based heuristics for natural gas
storage valuation'', \emph{Operations Research} 58(3), 2010.
\item US Energy Information Administration, Henry Hub natural gas spot price, daily, via FRED, consulté le 25 septembre 2026.
\end{omsources}

\section{Exercices}

\begin{exercise}[$\star$]\label{exo:s2:storage-transport-and-physical-optionality:1}
Calculer la valeur du plan intrinsèque~: 2 unités achetées chaque mois d'avril à août à 2.60~\$, 2.62~\$, 2.68~\$, 2.74~\$ et 2.78~\$~; 4 vendues à 3.45~\$ et
4 à 3.50~\$~; 2 à 3.30~\$~; 2 cents par MMBtu dans chaque sens~; unités de 100\,000 MMBtu.
\end{exercise}

\begin{exercise}[$\star$]\label{exo:s2:storage-transport-and-physical-optionality:2}
En janvier, le hub de la zone de marché se situe 60 cents au-dessus du hub et le transport coûte 10 cents. Que vaut une unité de capacité de janvier sur les contrats à terme~?
\end{exercise}

\begin{exercise}[$\star$]\label{exo:s2:storage-transport-and-physical-optionality:3}
Un chargement engagé vers l'Europe à 10.00~\$ peut être dirigé vers l'Asie moyennant 80 cents de fret supplémentaires. Quelle est la valeur intrinsèque de cette bascule lorsque l'Asie se situe à 10.50~\$,
et à 11.20~\$~?
\end{exercise}

\begin{exercise}[$\star\star$]\label{exo:s2:storage-transport-and-physical-optionality:4}
Pourquoi le programme glissant ne peut-il jamais rapporter moins que le programme statique sur un chemin quelconque lorsque le trading est gratuit~?
\end{exercise}

\begin{exercise}[$\star\star$]\label{exo:s2:storage-transport-and-physical-optionality:5}
Pourquoi réoptimiser un plan qui ignore l'incertitude capte-t-il l'essentiel de la valeur d'option~?
\end{exercise}

\begin{exercise}[$\star\star$]\label{exo:s2:storage-transport-and-physical-optionality:6}
Pourquoi la capacité de transport vaut-elle davantage lorsque les deux hubs sont moins corrélés~?
\end{exercise}

\begin{exercise}[$\star\star\star$]\label{exo:s2:storage-transport-and-physical-optionality:7}
\emph{Programmation.} Exécuter \texttt{gated(0.005)}. Comment le gating modifie-t-il le programme glissant à un demi-cent par unité~?
\end{exercise}

\begin{exercise}[$\star\star\star$]\label{exo:s2:storage-transport-and-physical-optionality:8}
\emph{Trouver la faille.} «~Le gaz d'hiver a été plus cher que le gaz d'été la plupart des années, donc nous louons du stockage, injectons tout l'été et vendons en hiver~; la couverture ne nous coûte que de l'argent.~»
ne fait que nous coûter de l'argent.~»
\end{exercise}

\section{Problème~: l'actif comme options}

\begin{problem}[{Problème du week-end --- stockage, transport et optionnalité physique}]\label{pb:s2:storage-transport-and-physical-optionality:1}
Le bail synthétique, le pipeline et le chargement du chapitre, ainsi que l'historique du Henry Hub.

\medskip\noindent\textbf{Partie I --- Stockage.}
\begin{enumerate}
\item Définir un \omterm{def:s2:storage-transport-and-physical-optionality:storage}{contrat de stockage}.
\item Décrire le bail, la courbe et le plan intrinsèque.
\item Donner les quatre valeurs et ce que suppose chaque méthode.
\item Que trouvent Secomandi, puis Lai, Margot et Secomandi~?
\end{enumerate}

\medskip\noindent\textbf{Partie II --- Monétisation.}
\begin{enumerate}[resume]
\item Décrire les programmes statique et glissant.
\item Pourquoi le programme statique rapporte-t-il la même chose sur chaque chemin~?
\item Donner les résultats du programme glissant pour chaque coût.
\item À quoi sert le gating~?
\end{enumerate}

\medskip\noindent\textbf{Partie III --- Transport et chargements.}
\begin{enumerate}[resume]
\item Définir une \omterm{def:s2:storage-transport-and-physical-optionality:transport}{option de transport}.
\item Valoriser le strip de transport, intrinsèquement et en tant qu'options.
\item Comment sa valeur dépend-elle de la corrélation~?
\item Valoriser la bascule de chargement.
\end{enumerate}

\medskip\noindent\textbf{Partie IV --- Le verdict.}
\begin{enumerate}[resume]
\item Énoncer le \emph{résultat nommé}~: la valeur du stockage selon la méthode et la valeur supplémentaire de la couverture glissante.
\item Que gagnait un opérateur non couvert au Henry Hub de 2007 à 2025~?
\item Définir le \omterm{def:s2:storage-transport-and-physical-optionality:assetbacked}{trading adossé à des actifs physiques}~; que possède le trader~?
\item Pourquoi la méthode de Monte-Carlo par moindres carrés est-elle inférieure à l'intrinsèque glissant ici~?
\item Pourquoi les paires d'options sur spread apportent-elles si peu ici~?
\item Combien paieriez-vous le bail~?
\item Quelle fiche de stratégie échouerait en premier si la courbe à terme devenait plate~?
\item En une phrase~: que vend un trader de stockage~?
\end{enumerate}
\end{problem}

\section{Questions d'entretien}

\begin{interviewq}[$\star$ \iqroles{trader}]\label{iq:s2:storage-transport-and-physical-optionality:1}
Quelle est la valeur intrinsèque d'une installation de stockage, et qu'est-ce que la valeur extrinsèque~?
\end{interviewq}

\begin{interviewq}[$\star\star$ \iqroles{chercheur}]\label{iq:s2:storage-transport-and-physical-optionality:2}
Comment valoriseriez-vous un stockage de gaz, et comment sauriez-vous que votre méthode est proche de l'optimum~?
\end{interviewq}

\begin{interviewq}[$\star\star$ \iqroles{trader}]\label{iq:s2:storage-transport-and-physical-optionality:3}
Le spread hiver-été s'élargit de 30 cents. Que faites-vous de vos couvertures de stockage~?
\end{interviewq}

\begin{interviewq}[$\star\star$ \iqroles{risque}]\label{iq:s2:storage-transport-and-physical-optionality:4}
Comment mesureriez-vous le risque d'un portefeuille de stockage entièrement couvert à l'intrinsèque~?
\end{interviewq}

\begin{interviewq}[$\star\star$ \iqroles{développeur}]\label{iq:s2:storage-transport-and-physical-optionality:5}
La réoptimisation de l'intrinsèque glissant doit être exécutée sur chaque chemin à chaque mois. Comment la rendriez-vous rapide~?
\end{interviewq}

\begin{interviewq}[$\star\star\star$ \iqroles{chercheur}]\label{iq:s2:storage-transport-and-physical-optionality:6}
Montrez que le P\&L du programme statique égale la valeur intrinsèque sur chaque chemin, quels que soient les cours au comptant, et dites quelles hypothèses cela requiert.
\end{interviewq}
